% ============================================================
% Section 120: 一维固体电子能谱与有效质量
% ============================================================
\section{固体物理：一维固体电子能谱与有效质量}
\label{sec:1d-band-structure}

\begin{modelbox}[题目：一维固体中电子的能谱分析]
图2.2为一维固体中电子的能谱。

(1) 若 $n$ 和 $p$ 分别为电子和空穴的数密度，则 $p/n$ 为何？

(2) 此固体每个原胞有偶数个还是奇数个电子？说明理由；

(3) 电子与空穴的有效质量，哪一个大？根据图示各量，求出有效质量的近似表示式。
\end{modelbox}

\begin{tipbox}[考点快速分析]
\begin{itemize}
    \item \textbf{一维 $k$ 空间态密度均匀}：粒子数正比于波矢区间长度
    \item \textbf{空穴定义}：满带中缺失的电子，$k_h=-k_e$
    \item \textbf{能带填充与奇偶性}：满带对应偶数个价电子（绝缘体/半导体），半满带对应奇数个价电子（金属）
    \item \textbf{有效质量}：$m^*=\hbar^2\left(\dd[2]{E}{k}\right)^{-1}$，曲率越大有效质量越小
\end{itemize}
\end{tipbox}

\begin{derivationbox}[标准解题过程]
\textbf{(1) 电子与空穴数密度之比 $p/n$}

由图2.2可见，在绝对零度下，电子填充至费米能级 $E_F$。$E_F$ 以下有一个被部分填充的能带：波矢区间 $(k_3,k_4)$ 被电子占据，而 $(k_1,k_2)$ 区间为空（导带底部）。根据空穴的定义，满带中缺失的电子即为空穴，故空穴占据的波矢区间为 $(k_1,k_2)$。

由于一维 $k$ 空间的态密度是均匀的，波矢区间长度正比于粒子数密度。由图可见两段区间长度相等，即
\[
k_2-k_1=k_4-k_3,
\]
因此电子和空穴的数密度相等：
\begin{equation}
\boxed{\frac{p}{n}=1.}
\end{equation}

\textbf{(2) 每个原胞电子数的奇偶性}

在绝对零度下，若不存在能带交叠，电子从最低能带开始填充。一维情况下每个能带可容纳 $2N$ 个电子（$N$ 为原胞数，因子 2 来自自旋简并）。图中 $E_F$ 以下刚好填满一个能带，说明总电子数为 $2N$，即每个原胞贡献 2 个电子。

因此，每个原胞含有 \textbf{偶数个电子}。理由：若每个原胞电子数为奇数，则能带必定部分填充，系统表现为金属；图中显示能带恰好填满，对应偶数个价电子，系统为半导体或绝缘体。

\textbf{(3) 有效质量的比较与近似表达式}

有效质量反比于能带曲率：
\[
m^*=\hbar^2\left(\dd[2]{E}{k}\right)^{-1}.
\]
曲率越大，有效质量越小。从图2.2可见：
\begin{itemize}
    \item 导带底 $(k_1,k_2)$ 区间曲线较陡，曲率大，故电子有效质量 $m_e^*$ 较小；
    \item 价带顶 $(k_3,k_4)$ 区间曲线较平缓，曲率小，故空穴有效质量 $m_h^*$ 较大。
\end{itemize}
结论：$\boxed{m_h^*>m_e^*}$。

\textit{近似表达式（抛物线近似）}：

对于导带电子，以导带底 $k_c=(k_1+k_2)/2$ 为原点：
\[
E(k)\approx E_c+\frac{\hbar^2(k-k_c)^2}{2m_e^*}.
\]
在 $k=k_2$ 处，$E(k_2)=E_F$，区间半宽为 $(k_2-k_1)/2$，故
\[
E_F-E_c=\frac{\hbar^2}{2m_e^*}\left(\frac{k_2-k_1}{2}\right)^2=\frac{\hbar^2(k_2-k_1)^2}{8m_e^*},
\]
解得
\begin{equation}
\boxed{m_e^*=\frac{\hbar^2(k_2-k_1)^2}{8(E_F-E_c)}.}
\end{equation}

对于价带空穴，以价带顶 $k_v=(k_3+k_4)/2$ 为原点，空穴能量向下为正：
\[
E(k)\approx E_v-\frac{\hbar^2(k-k_v)^2}{2m_h^*}.
\]
在 $k=k_4$ 处，$E(k_4)=E_F$，区间半宽为 $(k_4-k_3)/2$，故
\[
E_v-E_F=\frac{\hbar^2(k_4-k_3)^2}{8m_h^*},
\]
解得
\begin{equation}
\boxed{m_h^*=\frac{\hbar^2(k_4-k_3)^2}{8(E_v-E_F)}.}
\end{equation}
\end{derivationbox}

\begin{formulabox}[答案汇总]
\begin{center}
\begin{tabular}{ll}
\toprule
\textbf{问题} & \textbf{答案} \\
\midrule
(1) $p/n$ & 1（电子与空穴数密度相等） \\
(2) 原胞电子数 & 偶数个（能带恰好填满，对应二价） \\
(3) 有效质量大小 & $m_h^*>m_e^*$ \\
(3) 电子有效质量 & $m_e^*=\dfrac{\hbar^2(k_2-k_1)^2}{8(E_F-E_c)}$ \\
(3) 空穴有效质量 & $m_h^*=\dfrac{\hbar^2(k_4-k_3)^2}{8(E_v-E_F)}$ \\
\bottomrule
\end{tabular}
\end{center}
\end{formulabox}

\begin{insightbox}[物理图像]
\begin{itemize}
    \item 一维 $k$ 空间态密度均匀，粒子数正比于波矢区间长度。空穴与电子区间等长意味着本征激发下 $n=p$。
    \item 满带对应偶数个价电子，是绝缘体或半导体的特征；半满带对应奇数个价电子，是金属的特征。
    \item 有效质量由能带曲率决定：曲率越小（能带越平），有效质量越大。空穴通常比电子有效质量大，这是许多半导体的普遍规律。
\end{itemize}
\end{insightbox}
